\chapterbackmatter{Supplementary Exercises}

\begin{enumerate}
  \SupplementaryExercise{1}{Critical Scaling, \(O(N)\) Goldstone Modes, and a Phase Diagram}
  {Adapted from University of Cambridge, Part III Statistical Field Theory,
  Paper 303, Question 3 (2025).}
  {cambridge-part-iii-sft-2025-p303-q3-v2ch19-p01}
  \begin{enumerate}
    \item[(a)]
    Consider a critical point in the Ising universality class. Assume that,
    near the critical point, the most singular part of the free energy
    depends only on the correlation length. Derive a scaling law relating the
    critical exponents \(\alpha\) and \(\nu\).
    \item[(b)]
    Consider the \(O(N)\) model in \(d\) dimensions,
    \[
      F=\int d^dx\left[
        \frac12(\boldsymbol\nabla\boldsymbol\phi)^2
        +\frac12\mu^2\boldsymbol\phi^2
        +g(\boldsymbol\phi^2)^2+\cdots\right],
    \]
    where \(\boldsymbol\phi=(\phi_1,\ldots,\phi_N)\), \(N\geq2\), and
    \[
      \boldsymbol\phi^2=\sum_a\phi_a\phi_a,
      \qquad
      (\boldsymbol\nabla\boldsymbol\phi)^2
        =\sum_a\boldsymbol\nabla\phi_a\cdot
          \boldsymbol\nabla\phi_a.
    \]
    \begin{enumerate}
      \item[(i)]
      Use this model to explain the meaning of the terms spontaneous symmetry
      breaking and Goldstone boson.
      \item[(ii)]
      Use the quadratic approximation to show that Goldstone bosons have
      infinite correlation length.
      \item[(iii)]
      What is the lower critical dimension for this model? Explain your
      answer.
    \end{enumerate}
    \item[(c)]
    Consider a different theory with scalar fields
    \((\phi_1,\phi_2,\phi_3)\) and effective free energy
    \[
      F=\int d^dx\left[
        \frac12\sum_{i=1}^3(\boldsymbol\nabla\phi_i)^2
        +\frac12\sum_{i=1}^3\mu_i^2\phi_i^2
        +\sum_{i,j=1}^3g_{ij}\phi_i^2\phi_j^2\right],
    \]
    where \(g_{ij}\) is a real, symmetric, positive-definite matrix.
    \begin{enumerate}
      \item[(i)]
      The theory is invariant under a \(\mathbb Z_2\times O(2)\) symmetry,
      where \(\mathbb Z_2\) acts as \(\phi_1\mapsto-\phi_1\) and \(O(2)\)
      rotates \((\phi_2,\phi_3)\). What relations between the coupling
      constants does this symmetry imply? Explain your answer.
      \item[(ii)]
      Assume that
      \[
        g_{11}=g_{22}=g_{33}=g_{12}=g_{13}=g_{23}=g>0,
        \qquad \mu_3^2=\mu_2^2\neq\mu_1^2.
      \]
      Use mean-field theory to predict the phase diagram in the
      \((\mu_1^2,\mu_2^2)\) plane. What are the unbroken symmetries in each
      phase? Which phases have a Goldstone boson?
    \end{enumerate}
  \end{enumerate}

  \SupplementaryExercise{2}{Two-Flavor QCD and the Pion Effective Theory}
  {Adapted from University of Cambridge, Part III The Standard Model,
  Paper 305, Question 2 (2025).}
  {cambridge-part-iii-standard-model-2025-p305-q2-v2ch19-p02}
  Consider QCD coupled to two flavors of quarks, up and down, described by
  Dirac fermions \(q_i\), with \(i=1,2\). Take both quarks to be massless.
  \begin{enumerate}
    \item[(i)]
    Write down the Lagrangian for the theory. What is the global symmetry
    group of the classical theory? What is the global symmetry group of the
    quantum theory?
    \item[(ii)]
    Assume that a quark condensate forms, given by
    \[
      \langle\overline q_{L i}q_{R j}\rangle=-\sigma\delta_{ij},
    \]
    where \(q_L\) and \(q_R\) are left- and right-handed fermions,
    respectively, and \(\sigma\) is a dimensionful constant. Describe the
    resulting symmetry breaking. What is the vacuum manifold of this theory?
    Explain why this necessarily means that the low-energy physics, with no
    more than two derivatives, is described by the Lagrangian
    \[
      \mathcal L_{\mathrm{pion}}
        =\frac{f_\pi^2}{4}
          \operatorname{tr}(\partial^\mu U^\dagger\partial_\mu U)
    \]
    for a matrix-valued field \(U(x)\in SU(2)\). What is the dimension of the
    constant \(f_\pi\)?
    \item[(iii)]
    Write the field as
    \(U(x)=\exp(2i\pi(x)/f_\pi)\), with the pion field
    \(\pi(x)\in\mathfrak{su}(2)\), the Lie algebra of \(SU(2)\). Show that
    the terms quadratic and quartic in \(\pi\) take the form
    \[
      \mathcal L_{\mathrm{pion}}
        =\operatorname{tr}(\partial_\mu\pi)^2
        -\frac{2}{3f_\pi^2}
          \operatorname{tr}\!\left[
            \pi^2(\partial_\mu\pi)^2-(\pi\partial_\mu\pi)^2
          \right].
    \]
    \item[(iv)]
    The pion fields can be coupled to the \(U(1)\) of electromagnetism by
    introducing the charge matrix
    \[
      Q=\begin{pmatrix}\frac23&0\\[2pt]0&-\frac13\end{pmatrix}.
    \]
    The photon field \(A_\mu\) is coupled to \(U(x)\), with coupling constant
    \(e\), by promoting the derivative to a covariant derivative,
    \[
      D_\mu U=\partial_\mu U-ieA_\mu[Q,U].
    \]
    Explain, in terms of the symmetries of the theory, why this is the right
    coupling. By examining how the covariant derivative acts on the pion
    field, determine the electric charges of the various components of
    \(\pi\in\mathfrak{su}(2)\).
  \end{enumerate}

  \SupplementaryExercise{3}{An Abelian Coset with Charged Matter}
  {Adapted from Daniel Harlow, \emph{Relativistic Quantum Field Theory III},
  Section 9.8, Problem 3 (2026).}
  {harlow-qft3-2026-sec9p8-problem3-v2ch19-p03}
  Let
  \(G=U(1)_L\times U(1)_R\) be spontaneously broken to the diagonal
  subgroup
  \(H=\{(e^{i\beta},e^{i\beta})\}\).  Represent a point of \(G/H\) by
  \[
    g(x)=(e^{i\theta_L(x)},e^{i\theta_R(x)}),
    \qquad g\sim gh^{-1}.
  \]
  Decompose the Maurer--Cartan form orthogonally as
  \[
    g^{-1}\partial_\mu g=E_\mu+F_\mu,\qquad
    E_\mu\in\mathfrak h,\quad F_\mu\in\mathfrak h^\perp .
  \]
  \begin{enumerate}
    \item[(a)]
    Compute \(E_\mu\) and \(F_\mu\) in terms of
    \(\theta_L,\theta_R\).
    \item[(b)]
    Determine how \(\theta_L,\theta_R\) transform under a global element
    \((e^{i\alpha_L},e^{i\alpha_R})\in G\) and under the local change of
    representative \(g\mapsto gh^{-1}\).  Identify the physical Goldstone
    coordinate.
    \item[(c)]
    Let a complex field \(\Phi\) have charges \((q_L,q_R)\) under \(G\).
    Form the dressed field
    \[
      \widetilde\Phi
        =e^{-i(q_L\theta_L+q_R\theta_R)}\Phi .
    \]
    Find its transformation under the local \(H\) redundancy and construct
    \(D_\mu\widetilde\Phi\) using \(E_\mu\).  Write the result explicitly
    in terms of \(\theta_L,\theta_R,\Phi\).
    \item[(d)]
    Express
    \(\operatorname{Tr}(F_\mu F^\mu)\) and
    \((D_\mu\widetilde\Phi)^\dagger
      D^\mu\widetilde\Phi\)
    in the original variables, and verify invariance under both the global
    \(G\) action and the local choice of coset representative.
  \end{enumerate}

  \SupplementaryExercise{4}{Goldstone Power Counting, Curvature, and Pseudo-Goldstones}
  {Inspired by David Tong, Statistical Field Theory (2021).}
  {tong-statistical-field-theory-goldstone-2021-v2ch19-p04}
  \begin{enumerate}
    \item[(a)]
    A connected graph has \(L\) loops and \(V_i\) vertices containing
    \(d_i\) derivatives.  Derive
    \(\nu=2+2L+\sum_iV_i(d_i-2)\) for its low-momentum order \(Q^\nu\).
    Apply the formula to a one-loop graph built only from two-derivative
    vertices.
    \item[(b)]
    In Riemann normal coordinates about the vacuum,
    \(g_{ab}=\delta_{ab}-R_{acbd}\pi_c\pi_d/3+\cdots\).
    Show that the leading four-Goldstone interaction is fixed by the target
    curvature.  Specialize to a sphere of radius \(F\) and state the
    crossing-symmetric massless amplitude.
    \item[(c)]
    An \(O(3)\) order parameter of fixed length \(F\) is weakly perturbed by
    \(\Delta\mathcal L=hF n_3\).  Expand about the aligned vacuum and find
    the two pseudo-Goldstone masses.  Explain the dependence on \(h\).
  \end{enumerate}

  \SupplementaryExercise{5}{Current Poles, PCAC, and Pion Thresholds}
  {Inspired by Stefan Scherer, Introduction to Chiral Perturbation Theory (2002).}
  {scherer-chiral-perturbation-2002-v2ch19-p05}
  \begin{enumerate}
    \item[(a)]
    Let \(J^\mu_a\) be a broken conserved current and \(\Phi_b\) an operator
    with a nonzero vacuum commutator.  Use the spectral representation of
    \(\langle T J^\mu_a(x)\Phi_b(0)\rangle\) and current conservation to
    show that its Fourier transform must contain a pole at \(q^2=0\).
    \item[(b)]
    Given
    \(\bra{\Omega}A^\mu_a(0)\ket{\pi_b(p)}
      =if_\pi p^\mu\delta_{ab}\),
    derive the one-pion matrix element of \(\partial_\mu A^\mu_a\) for a
    pion of mass \(m_\pi\).  State what happens in the chiral limit.
    \item[(c)]
    Starting from
    \(A(s,t,u)=(s-m_\pi^2)/f_\pi^2\), form the \(I=0,1,2\)
    amplitudes
    \(T^0=3A(s,t,u)+A(t,s,u)+A(u,t,s)\),
    \(T^1=A(t,s,u)-A(u,t,s)\), and
    \(T^2=A(t,s,u)+A(u,t,s)\).
    Evaluate them at threshold and obtain \(a_0^0\) and \(a_0^2\) in the
    convention \(a_0^I=T^I_{\rm thr}/(32\pi)\).
  \end{enumerate}

  \SupplementaryExercise{6}{Goldberger--Treiman, Baryon Counting, and Soft-Pion Commutators}
  {Inspired by Stefan Scherer, Introduction to Chiral Perturbation Theory (2002).}
  {scherer-chiral-perturbation-2002-v2ch19-p06}
  \begin{enumerate}
    \item[(a)]
    The leading pion--nucleon interaction is
    \(-g_A\overline N\gamma^\mu\gamma_5\tau_a
    (\partial_\mu\pi_a)N/(2f_\pi)\).
    Use the on-shell Dirac equations to replace it by a pseudoscalar
    Yukawa interaction and derive the Goldberger--Treiman relation.
    \item[(b)]
    Count a pion propagator as \(Q^{-2}\), a nearly on-shell nucleon
    propagator as \(Q^{-1}\), a loop as \(Q^4\), and a vertex of type \(i\)
    as \(Q^{d_i}\).  Derive the chiral index for a connected graph with
    \(E_N\) external nucleon lines and vertices containing \(n_i\) nucleon
    fields.
    \item[(c)]
    For a local operator \(\mathcal O\), use the pion pole in the axial
    current to show schematically that
    \(\lim_{q\to0}\bra{\beta;\pi^a(q)}\mathcal O\ket{\alpha}\)
    is fixed by \(\bra{\beta}[Q_5^a,\mathcal O]\ket{\alpha}/f_\pi\),
    apart from external-leg pole terms.  Track the factor of \(i\).
  \end{enumerate}

  \SupplementaryExercise{7}{Running Couplings, Low-Energy QCD, and Chiral Symmetry}
  {Adapted from University of Cambridge, Part III The Standard Model,
  Paper 305, Question 4 (2023).}
  {cambridge-part-iii-standard-model-2023-p305-q4-v2ch19-p07}
  \begin{enumerate}
    \item[(a)]
    The scale dependence of a renormalized coupling \(g(\mu)\) is given, to
    leading order, by the beta-function equation
    \[
      \mu\frac{d\alpha}{d\mu}=-\frac{b}{2\pi}\alpha^2,
    \]
    where \(\mu\) is the renormalization scale and \(b\) is a real constant.
    Consider an \(SU(N_c)\) gauge theory with \(N_f\) flavors in the
    fundamental representation of \(SU(N_c)\), for which
    \[
      b=\frac13(11N_c-2N_f).
    \]
    For QCD, \(N_c=3\), and the gauge coupling will be denoted by \(g_3\).
    \begin{enumerate}
      \item[(i)]
      Solve the beta-function equation for \(\alpha_3=g_3^2/(4\pi)\). Sketch
      the profile of \(\alpha_3(\mu)\), and estimate the scale
      \(\Lambda_{\mathrm{QCD}}\), at which QCD becomes strongly coupled, as
      a function of \(b\) and the renormalization scale \(\mu\).
      \item[(ii)]
      Briefly explain the difference between theories in which \(b<0\) and
      \(b>0\).
      \item[(iii)]
      Suppose there is an energy scale \(M_{\mathrm{GUT}}>M_Z\), where
      \(M_Z\) is the mass of the \(Z\) boson, at which the three gauge
      couplings of the Standard Model coincide,
      \[
        \alpha_1(M_{\mathrm{GUT}})=\alpha_2(M_{\mathrm{GUT}})
          =\alpha_3(M_{\mathrm{GUT}}),
      \]
      corresponding to \(U(1)\), \(SU(2)_L\), and \(SU(3)_c\), respectively.
      Show that this implies that at the \(M_Z\) scale they are related by
      \[
        \frac1{\alpha_3(M_Z)}
        =\frac1{\alpha_2(M_Z)}
        +\frac{b_3-b_2}{b_1-b_2}
          \left[\frac1{\alpha_1(M_Z)}-\frac1{\alpha_2(M_Z)}\right],
      \]
      where \(b_i\) are the beta-function coefficients for the couplings
      \(\alpha_i\), respectively.
    \end{enumerate}
    \item[(b)]
    Keeping in mind that the masses of the quarks are known to be
    \(m_{u,d}\sim O(1)\,\mathrm{MeV}\), with the other ones much heavier,
    \[
      m_s,m_c,m_b,m_t\sim
        10^2,10^3,4\times10^3,1.7\times10^5\ \mathrm{MeV},
    \]
    respectively, write down the low-energy QCD Lagrangian including only
    the \(u\) and \(d\) quarks, consistent with the gauge symmetry
    \(SU(3)_c\times U(1)_{\mathrm{EM}}\). Show that baryon number is an
    accidental symmetry of this Lagrangian.
    \item[(c)]
    Identify the approximate chiral symmetries of this Lagrangian in the
    limit \(m_{u,d}\to0\), and argue why they can be broken.
    \begin{enumerate}
      \item[(i)]
      Construct an effective field theory for scalar fields that captures
      the desired symmetry breaking. Identify the unbroken symmetry group
      and the identity of the pseudo-Goldstone modes of this approximate
      symmetry.
      \item[(ii)]
      Write down an effective Lagrangian for the pseudo-Goldstone bosons for
      energies \(E\ll\Lambda_{\mathrm{QCD}}\).
      \item[(iii)]
      Could this approximate symmetry explain the mass difference between
      protons and neutrons?
    \end{enumerate}
    \item[(d)]
    Is it possible to extend this analysis to include also the \(s\) and
    \(c\) quarks? Explain.
  \end{enumerate}

  \SupplementaryExercise{8}{Quark-Mass Ratios, Loop Counterterms, and the WZW Term}
  {Inspired by Antonio Pich, Effective Field Theory (1998).}
  {pich-effective-field-theory-1998-v2ch19-p08}
  \begin{enumerate}
    \item[(a)]
    Assume the leading electromagnetic shift is the same for
    \(K^+\) and \(\pi^+\), and vanishes for \(K^0,\pi^0\).
    Eliminate the common electromagnetic term and \(B_0\) to derive
    \(m_u/m_d\) from the four measured squared meson masses.
    \item[(b)]
    Use chiral power counting to show that one-loop graphs made only from
    \(\mathcal L_2\) and tree graphs with one \(\mathcal L_4\) insertion
    both contribute at order \(Q^4\).  Explain why the coefficients in
    \(\mathcal L_4\) must run with the subtraction scale.
    \item[(c)]
    Put \(U=\exp(2iB/F)\) in
    \(\int_{B_5}\operatorname{Tr}(U^{-1}dU)^5\).
    Determine the lowest power of \(B\), the number of derivatives, and the
    parity-odd tensor that appears in four-dimensional spacetime.  Explain
    why its coefficient is quantized.
  \end{enumerate}

  \SupplementaryExercise{9}{Matrix-Scalar Vacua and Goldstone Bosons}
  {Adapted from Daniel Harlow, Physics 8.324,
  \emph{Relativistic Quantum Field Theory II}, Section 13.7, Problem 3
  (Fall 2024).}
  {harlow-qft2-2024-sec13p7-problem3-v2ch19-p09}
  Consider a scalar field theory in which \(\phi\) is an \(N\times N\)
  matrix of scalar fields and the Lagrangian is
  \[
    \mathcal L
      =-\operatorname{Tr}(\partial_\mu\phi^\dagger\partial^\mu\phi)
       +m^2\operatorname{Tr}(\phi^\dagger\phi)
       -\frac{\lambda}{4}
        \operatorname{Tr}(\phi^\dagger\phi\phi^\dagger\phi).
  \]
  There is an \(SU(N)\) global symmetry that acts on the field as
  \[
    \phi' = U^\dagger\phi U.
  \]
  Strictly speaking, the symmetry group is \(SU(N)/\mathbb Z_N\), since the
  central elements of \(SU(N)\) act trivially on \(\phi\), but this will not
  be important in this problem.
  \begin{enumerate}
    \item[(a)]
    What is the set of minima of the potential for \(\phi\)? Hint: think
    about diagonalizing \(\phi^\dagger\phi\).
    \item[(b)]
    For a generic minimum \(\phi_0\), what subgroup leaves \(\phi_0\)
    invariant?
    \item[(c)]
    How many Goldstone bosons does the theory have for a generic \(\phi_0\)?
    Can you think of a special \(\phi_0\) for which it has fewer?
  \end{enumerate}

  \SupplementaryExercise{10}{Positivity, \(U(1)_A\), and the Neutral-Pion Anomaly}
  {Inspired by MIT OpenCourseWare 8.325, Relativistic Quantum Field Theory III (2007).}
  {mit-ocw-qft-iii-2007-v2ch19-p10}
  \begin{enumerate}
    \item[(a)]
    Let \(\slashed D\) be anti-Hermitian in Euclidean space and let all
    eigenvalues of the Hermitian flavor mass matrix satisfy
    \(|m_a|\geq m_{\min}>0\).  Prove
    \(\|(\slashed D+M)^{-1}\|\leq1/m_{\min}\), and explain why this blocks
    spontaneous breaking of exact vector flavor symmetries by a singular
    zero-mass limit.
    \item[(b)]
    Write the classical singlet axial current of \(N_f\) massless quarks and
    its quantum divergence in a non-Abelian gauge field.  Explain why the
    anomaly removes the ninth Goldstone prediction for \(N_f=3\).
    \item[(c)]
    Using
    \(\mathcal L_{\pi\gamma\gamma}
      =g_{\pi\gamma\gamma}\pi^0
        \epsilon^{\mu\nu\rho\sigma}F_{\mu\nu}F_{\rho\sigma}/8\)
    with \(g_{\pi\gamma\gamma}=\alpha/(\pi f_\pi)\), derive
    \(\Gamma(\pi^0\to\gamma\gamma)\) in terms of
    \(\alpha,f_\pi,m_\pi\).
  \end{enumerate}

  \SupplementaryExercise{11}{The Two-Dimensional Goldstone Obstruction}
  {From John McGreevy, Physics 215C, Assignment 4, Problem 3 (2025), with
  typographical normalization.}
  {mcgreevy-215c-s25-a4-p3-v2ch19-p11}
  In this problem we study a free massless boson in two dimensions.  This
  solvable system has many physical applications---for example in string
  theory and at the edge of quantum Hall systems.  It is an example of a
  conformal field theory, of a field theory whose excitations are not
  particles, and of a theory in which all the composite operators can be
  understood exactly.

  \begin{enumerate}
    \item[(a)] Consider a massless scalar \(X\) in two dimensions, with action
    \[
      S[X]= -\frac{1}{4\pi}\int d^2\sigma\,
      \partial_aX\,\partial^aX .
      \tag{1}
    \]
    Show that the Euclidean Green function \(G_2\) satisfies
    \[
      \nabla^2G_2(z,z')=-2\pi\delta^2(z-z')
      \tag{2}
    \]
    and, writing \(z=\sigma_E^1+i\sigma_E^2\), is given by
    \[
      G_2(z,z')=-\ln|z-z'|,
    \]
    for example by Fourier transform.

    \emph{Hint.} Recall the Schwinger--Dyson equations.

    \item[(b)] The long-distance behavior of \(G_2\) has important implications
    for the physics of massless scalars in two dimensions.  Thinking of
    \(G_2\) as the two-point function of the massless scalar,
    \[
      G_2(z,z')=\langle X(z)X(z')\rangle,
    \]
    let us ask the following question.

    There is no potential energy for the field \(X\) in (1).  Someone used
    to doing physics in \(3+1\) dimensions might think that this means there
    is a vacuum for every value of \(X\).  Try to fix the expectation value
    of the scalar, \(\langle X\rangle=x\), and consider what happens.  Perturb
    the putative vacuum \(\ket{x}\) a little at the position \(z\) by
    inserting the operator \(X\) there.  To measure what happens, insert the
    operator \(X\) at \(z'\).  The correlator \(G_2\) can thus be interpreted
    as a measurement of how the effects of the perturbation fall off with
    distance.  What happens?  Contrast this with the behavior for a scalar
    field with a flat potential in more than two dimensions.  Note that the
    case of \((0+1)\)-dimensional QFT---that is, quantum mechanics---is even
    more problematic in the infrared.

    One way to arrive at an action such as (1) is for \(X\) to arise as a
    Goldstone boson associated with a symmetry \(X\mapsto X+a\), which would
    be broken by fixing the vacuum \(\ket{x}\).  Goldstone's theorem
    then guarantees that the action can depend only on derivatives of
    \(X\).  Explain the implication of the behavior found above for such a
    putative Goldstone boson.  The Goldstone nature of the massless boson is
    not crucial for the infrared discussion: massless bosons not protected
    by a symmetry would normally be expected to acquire a mass quantum
    mechanically, but special cases such as supersymmetric theories can
    possess distinct minima that are not related by a symmetry.

    This result is called the Coleman--Mermin--Wagner (and sometimes
    Hohenberg) theorem; see S.~Coleman, ``There are no Goldstone bosons in
    two dimensions,'' \emph{Commun. Math. Phys.} \textbf{31}, 259--264
    (1973).
  \end{enumerate}

\end{enumerate}
